\begin{problem}[第32届全国中学生物理竞赛决赛]{54}
飞行时间质谱仪(TOFMS)的基本原理如图1所示，主要由离子源区、漂移区和探测器三部分组成。带正电的离子在离子源中形成后被电场 $E_{s}$ 加速，经过漂移区（真空无场），到达离子探测器。设离子在离子源区加速的距离为S，在漂移区漂移的距离为L。通过记录离子到达探测器的时间，可以把不同的离子按质荷比m/q的大小进行分离，这里m和q分别表示离子的质量和电量。分辨率是TOFMS最重要的性能指标，本题将在不同情况下进行讨论或计算。忽略重力的影响。
\begin{subproblem}
对于理想的 TOFMS，不同离子在离子源 X 轴方向同一位置、同一时刻 t=0 产生，且初速度为 0。探测器可以测定离子到达探测器的时刻 t，其最小分辨时间为 $\Delta t$（即探测器所测时刻 t 的误差）。定义仪器的分辨率为 $R=m/\Delta m$，其中 $\Delta m$ 为最小分辨时间 $\Delta t$ 对应的最小分辨质量。此种情形下，R 完全由 $\Delta t$ 决定，试推导 R 与 t 和 $\Delta t$ 之间的关系。
\end{subproblem}
\begin{subproblem}
实际上，离子产生的位置也有微小的差别 $\Delta S$（$\Delta S \ll S$），这导致具有相同质荷比的离子不能同时到达探测器，从而影响质谱仪的分辨率。如图2所示，在离子源后引入第二加速电场 $E_{1}$，该电场区域的长度为 $D_{1}$，通过适当选择漂移区的长度 $L$，可使同一时刻在不同位置产生的质荷比相同的离子尽量同时到达探测器，以使分辨率受 $\Delta S$ 的影响最小，试求 $L$ 的值。
\end{subproblem}
\begin{subproblem}
为进一步降低离子产生位置的离散性对分辨率的影响，通常采用如图3所示的反射式TOFMS。这里在二级加速电场（$E_{s}$ 和 $E_{1}$）的基础上增加了反射器，它由两级电场 $E_{2}$ 和 $E_{3}$ 组成，通过这两级电场对离子的飞行方向进行反转，以使分辨率受 $\Delta S$ 的影响最小，试求 L 的值。为简化计算，假设离子的运动是平行于 X 方向的直线运动。（装置的各参数间满足关系 $E_{s}S + E_{1}D_{1} \leq E_{2}D_{2} + E_{3}D_{3}$，以使所有离子飞行方向的反转都可以实现。）
\end{subproblem}
\end{problem}